\documentclass[11pt,a4paper]{article}
\usepackage{fontspec}
\setmainfont{DejaVu Sans}
\usepackage[french]{babel}
\usepackage[margin=2.2cm]{geometry}
\usepackage{microtype}
\usepackage{hyperref}
\hypersetup{colorlinks=true,linkcolor=blue}
\emergencystretch=5em
\sloppy
\setlength{\parindent}{0pt}
\setlength{\parskip}{4pt}
\begin{document}
\begin{center}
{\LARGE Étude de faisabilité openEMS}\par
{\large LiteVNA 64, antenne UWB Surveyor et plaque métallique}\par
17 octobre 2026
\end{center}
\tableofcontents\newpage
\section{Question initiale de l’utilisateur}
Est-ce que OpenEMS saurait simuler :\par
\par\medskip
- un liteVNA \par
- connecté à l'antenne dont la datasheet est en pièce jointe ("uwb\_surveyor\_datasheet\_revision\_1.pdf")?\par
- avec une plaque de dimension 60x60cm\textasciicircum{}2 en métal placé à une distance de x cm de l'antenne et parallèle au plan de l'antenne \par
\par\medskip
afin de produire des figures comme celles incluses dans le document "report\_john.pdf" ci joint ?\par
\newpage
\section{Réponse précédente de l’assistant}
Oui, openEMS est capable de simuler ce montage complet et de produire les mêmes types de figures que celles de report\_john.pdf. C'est même une bonne approche pour étudier l'origine du décalage de distance d'environ 20 cm observé dans tes mesures.\par
\par\medskip
J'ai examiné les deux PDF. Le point essentiel est qu'openEMS peut calculer la réponse électromagnétique du système, mais qu'il faudra modéliser suffisamment fidèlement l'antenne pour obtenir des résultats quantitativement comparables aux mesures.\par
\par\medskip
\subsection*{1. Le montage à simuler}
\par\medskip
Vue de profil schématique : LiteVNA 64 (port 50 Ω) — câble coaxial — antenne UWB Surveyor — air (distance x cm) — plaque métallique 60 × 60 cm, parallèle au plan de l'antenne. Dimensions non à l'échelle.\par
\par\medskip
Les éléments seraient les suivants :\par
\par\medskip
| Élément | Modélisation openEMS |\par
| LiteVNA 64 | Port RF 50 Ω, sans simuler l'électronique interne |\par
| Câble coaxial | Ligne de transmission ou géométrie 3D |\par
| Antenne UWB Surveyor | Modèle électromagnétique 3D |\par
| Plaque 60 × 60 cm | Conducteur parfait (PEC) en première approximation |\par
| Distance x | Paramètre variable |\par
| Environnement | Air avec frontières absorbantes PML |\par
\par\medskip
\subsection*{2. Reproduire exactement le traitement de ton rapport}
\par\medskip
Ton rapport utilise les paramètres suivants :\par
\par\medskip
| Paramètre | Valeur |\par
| Bande | 1–6 GHz |\par
| Échantillons fréquentiels | 201 |\par
| Pas fréquentiel | 25 MHz |\par
| Distances étudiées | 50, 60, 70 cm |\par
| Référence | 50 Ω |\par
| Fenêtre | Hann symétrique |\par
| IFFT | 3216 points |\par
| Calibration expérimentale | SOL à l'extrémité du câble |\par
\par\medskip
Pour reproduire la figure 2, il faudra réaliser quatre simulations : sans plaque, puis avec la plaque à 50, 60 et 70 cm.\par
\par\medskip
On obtiendra :\par
\par\medskip
S11,background(f)\par
\par\medskip
et\par
\par\medskip
S11,x(f)\par
\par\medskip
On calculera ensuite, exactement comme dans ton rapport :\par
\par\medskip
R\_x(f) = S11,x(f) − S11,background(f)\par
\par\medskip
Puis :\par
\par\medskip
P\_x = 16 |IFFT\_3216(w(f) R\_x(f))|\par
\par\medskip
avec la conversion en distance :\par
\par\medskip
d\_m = c m / (2 × 3216 × 25 × 10\textasciicircum{}6)\par
\par\medskip
Cela permettra de produire les mêmes courbes distance/amplitude que dans la figure 2, avec le même traitement numérique. La normalisation des amplitudes simulées devra être cohérente avec celle des paramètres S mesurés.\par
\par\medskip
\subsection*{3. Le point scientifique le plus intéressant : le décalage de distance}
\par\medskip
Ton rapport montre deux branches de pics :\par
\par\medskip
| Distance réelle | Premier pic | Second pic |\par
| 50 cm | 69,91 cm | 75,88 cm |\par
| 60 cm | 79,61 cm | 86,32 cm |\par
| 70 cm | 90,42 cm | 96,57 cm |\par
\par\medskip
Le premier pic présente un décalage moyen d'environ 20 cm, et le second d'environ 26,3 cm.\par
\par\medskip
Avec openEMS, nous pourrions chercher à déterminer si ces décalages proviennent de l'antenne elle-même.\par
\par\medskip
Il serait notamment possible de calculer la phase complexe de la réponse à la plaque après suppression du trajet géométrique :\par
\par\medskip
H\_x(f) = R\_x(f) exp(j 4πfx/c)\par
\par\medskip
Puis d'en déduire un retard de groupe apparent :\par
\par\medskip
τ\_g(f) = −(1/2π) d(arg H\_x(f))/df\par
\par\medskip
On pourrait ainsi comparer les retards simulés à ceux du rapport, qui correspondent à des décalages équivalents de 21,3 à 22,7 cm.\par
\par\medskip
C'est probablement l'apport principal de la simulation : distinguer un retard introduit par la structure rayonnante de l'antenne d'un effet provenant du câble, des supports ou des réflexions dans la pièce.\par
\par\medskip
\subsection*{4. La principale difficulté : reproduire l'antenne}
\par\medskip
La datasheet indique un PCB de 120 × 120 × 0,8 mm, mais elle ne fournit pas les fichiers de fabrication ni la géométrie exacte des conducteurs.\par
\par\medskip
Je proposerais donc une simulation progressive :\par
\par\medskip
1. Modèle initial : antenne biconique planaire approximative, avec alimentation 50 Ω et plaque métallique.\par
2. Modèle affiné : géométrie des conducteurs, substrat PCB et structure d'alimentation plus réalistes.\par
3. Comparaison : ajustement du modèle à partir des courbes S11 mesurées sans plaque.\par
4. Validation : comparaison des profils de distance et des phases pour les trois positions de plaque.\par
\par\medskip
Il faudra éviter d'ajuster arbitrairement un retard de 1,4 ns simplement pour reproduire les 21 cm observés : l'objectif est précisément d'en identifier l'origine physique.\par
\par\medskip
\subsection*{5. Est-ce réalisable en Python ?}
\par\medskip
Oui. Je recommande une architecture Python avec openEMS pour les simulations FDTD, NumPy/SciPy pour le traitement fréquentiel et temporel, et Matplotlib pour les figures.\par
\par\medskip
Nous pourrions produire automatiquement :\par
\par\medskip
- Les fichiers Touchstone .s1p simulés.\par
- Les courbes S11 complexes.\par
- Les profils IFFT en fonction de la distance.\par
- Les deux branches de pics et leurs décalages.\par
- Les courbes de phase corrigées du trajet géométrique.\par
- Les champs E et H autour de l'antenne et de la plaque.\par
\par\medskip
Une réserve pratique : une simulation FDTD 3D couvrant une plaque de 60 cm et une fréquence maximale de 6 GHz peut être coûteuse en mémoire et en temps de calcul. Un maillage adaptatif spatialement et plusieurs niveaux de précision seront nécessaires.\par
\par\medskip
\subsection*{Ma proposition}
\par\medskip
Je commencerais par un projet Python/openEMS paramétrable, avec la distance x\_cm, la fréquence, la géométrie de l'antenne et les paramètres de maillage modifiables, capable de générer les figures 2 et 3 de ton rapport.\par
\par\medskip
Il y a toutefois une question importante avant de choisir le modèle :\par
\par\medskip
Souhaites-tu que l'on cherche à reproduire quantitativement le décalage d'environ 20 cm de tes mesures réelles, ou simplement à simuler le principe physique du montage ?\par
\par\medskip
Dans le premier cas, il faudra porter une attention particulière à la géométrie exacte de l'antenne, à sa structure d'alimentation et au plan de référence de la calibration SOL.\par
\newpage
\section{Nouveau prompt de l’utilisateur}
Je voudrais commencer par un projet Python/openEMS paramétrable qui permettrait :\par
\par\medskip
(1) simplement à simuler le principe physique du montage\par
\par\medskip
(2) cherche à reproduire quantitativement le décalage d'environ 20 cm de nos mesures réelles\par
\par\medskip
Est-ce possible ? Ne le fais pas, dis moi juste si c'est possbible et produit un document latex avec ma question "Est-ce que OpenEMS saurait simuler :\par
\par\medskip
- un liteVNA \par
- connecté à l'antenne dont la datasheet est en pièce jointe ("uwb\_surveyor\_datasheet\_revision\_1.pdf")?\par
- avec une plaque de dimension 60x60cm\textasciicircum{}2 en métal placé à une distance de x cm de l'antenne et parallèle au plan de l'antenne \par
\par\medskip
afin de produire des figures comme celles incluses dans le document "report\_john.pdf" ci joint ?"\par
\par\medskip
Ta réponse ci-dessus,\par
\par\medskip
Ce prompt que je suis en train d'écrire et la réponse que tu vas produire.\par
\par\medskip
COmpile le latex en pdf également.\par
\newpage
\section{Réponse à ce nouveau prompt}
Oui, les deux objectifs sont réalisables dans un même projet Python/openEMS paramétrable, avec deux niveaux de modélisation et de validation.\par
\par\medskip
Le premier objectif, simuler le principe physique du montage, est réalisable avec une antenne biconique planaire approximative, un port RF de 50 Ω représentant le LiteVNA 64 et une plaque métallique de 60 × 60 cm placée à une distance variable x. Il permettra de calculer S11(f), les réflexions, les profils IFFT et les phases, puis de produire des figures comparables à celles de report\_john.pdf.\par
\par\medskip
Le deuxième objectif, reproduire quantitativement le décalage apparent d'environ 20 cm, est également envisageable, mais sa réussite n'est pas garantie avec les seules informations disponibles. Il faudra affiner la géométrie de l'antenne, sa ligne d'alimentation, le plan de référence de la calibration et éventuellement les éléments de l'environnement expérimental. Il faudra aussi comparer les résultats aux données complexes S11 réellement mesurées.\par
\par\medskip
Je recommande donc un projet unique comportant deux étapes successives : un modèle physique simplifié, puis un modèle affiné et confronté aux mesures. Cette démarche permettra de distinguer les phénomènes expliqués par la simulation des contributions encore inconnues.\par
\par\medskip
Aucun développement ni aucune simulation openEMS n'a été lancé dans le cadre de cette demande. Ce document rassemble les échanges demandés et est fourni en source LaTeX et en PDF compilé.\par
\newpage
\end{document}
